\section{Discussion}

We present the first defense against cross-model adversarial transfer in LLMs, reducing cross-model ASR to 0--5\% across seven models spanning five architecture families where existing defenses leave 1.2--6.2\% residual ASR even in the same-model setting. The approach is grounded in a simple insight: representation similarity enables adversarial transfer, so breaking that similarity should prevent it. By minimizing CKA~\citep{kornblith2019similarity} between a defender and an anchor model, we disrupt the shared geometric structure that transferred attacks exploit. Because CKA operates on Gram matrices, the defense is dimension-agnostic and attack-agnostic---no adversarial examples are needed during training.

The ablation study reveals an important design lesson. CKA repulsion alone is sufficient for security (Run~2 achieves 0\% ASR) but collapses benign output quality (BGR $>$ 50\%). This failure mode---degenerate output that passes standard refusal-keyword evaluation---motivates our proposal of the Benign Garble Rate (BGR) metric, which we believe is essential for honest assessment of any representation-engineering defense. The full defense resolves this tension by combining moderate CKA repulsion with refusal-direction steering and KL-divergence preservation, achieving near-zero cross-model ASR with 0\% BGR.

The layer-depth analysis further connects diagnostic observations to defense design. Per-layer CKA diagnostics reveal that harmful-prompt similarity is uniformly high across all layers, while benign-prompt similarity is low at early layers and plateaus at mid-depth. This predicts the optimal defense operating point: applying CKA repulsion at mid-depth (layer 0.5) selectively disrupts harmful representations---where both harmful and benign CKA are high---without collapsing the model's general capabilities. Earlier layers destroy output quality despite achieving perfect security, while later layers preserve quality but fail to fully block transferred attacks.

\paragraph{Deployment overhead.}
The defense produces a standard LoRA adapter; after merging, inference latency and memory are identical to the base model~\citep{lora}. Training completes in 200--300 steps (${\sim}$15 minutes on a single L40S GPU).

\paragraph{Reproducibility.}
Evaluation uses greedy decoding over 2,000 attack prompts per defender (100 prompts $\times$ 20 source models), and the consistent 0--5\% cross-model ASR across seven independently trained models provides cross-model replication.

Several limitations warrant discussion. First, our defense targets the cross-model transfer channel specifically; it does not improve robustness against same-model adaptive attacks, and embedding-space PGD~\citep{madry2018towards} bypasses all representation-engineering defenses including ours~\citep{revisiting_cb}. More subtly, an attacker who knows both the defender and anchor identity could in principle optimize a suffix that avoids the anchor's trajectory while still targeting the defender---a defense-aware adaptive attack. Our current evaluation does not test this scenario, and we expect it would partially degrade protection, consistent with the broader adversarial arms race in this subfield. Second, the defense requires selecting an anchor model from a different family, and the optimal hyperparameters (particularly the CKA weight $\gamma$) vary across defender--anchor pairs. Third, while we demonstrate our defense across seven models spanning five architecture families, our evaluation covers the 7B--14B parameter scale. Because our multi-objective loss requires extracting intermediate hidden states and computing CKA Gram matrices across large batch sizes to ensure stable representations, scaling this specific training pipeline to frontier models (e.g., 70B+ parameters) introduces significant memory overhead. Given that 7B--14B models are currently the most accessible white-box targets for adversarial transfer, we prioritized this high-risk scale, leaving the engineering optimizations required for larger models to future work.

Looking forward, the representation-similarity framework opens several promising directions. Diagnostically, quantifying representation overlap could inform model-release decisions by flagging high-transfer-risk pairs before deployment. The defense mechanism itself could be extended to multi-anchor settings, repelling from several reference models simultaneously to provide broader protection. Finally, investigating whether CKA repulsion interacts synergistically with standard safety training methods (e.g., RLHF, DPO) could yield models that are intrinsically robust to both same-model and cross-model attacks.
